\scanpage{264}
\mspage{133 quinquies}

\[
 g_0\alpha_{j_1}=g_1\alpha_{j_1}
\]

i.e. les deux sections de $\unsure{\varepsilon|E}$ sur $e$ sont égales
dans $\mathrm{Hot}_{W}$. Comme le foncteur $\mathrm{Cat}\unsure{-|\to}\mathrm{Hot}_{W}$
commute aux produits finis, il s'ensuit
que \inserted{pour $X$ dans $\mathrm{Cat}$,} les deux flèches

\diagram{\begin{tikzcd}[ampersand replacement=\&,column sep=large]
X \arrow[r,shift left=0.65ex,"{\unsure{\partial|\delta}}_0^X"] \arrow[r,shift right=0.65ex,"\partial_1^X"'] \& X\times\varepsilon=X\amalg X
\end{tikzcd}}

sont égales dans $\mathrm{Hot}_{W}$. Si alors

\[
 f_0,f_1:X\mathrel{\unsure{=|\rightrightarrows}}Y
\]

sont deux flèches dans $\mathrm{Cat}$, elles sont de
la forme

\[
 f_0=h\partial_0^X,\qquad f_1=h\partial_1^X,
\]

avec

\[
 h=(f_0,f_1):X\amalg X\longrightarrow\unsure{X|Y},
\]

et comme $\partial_0^{\unsure{\alpha|X}}=\partial_1^X$ dans $\mathrm{Hot}_{W}$, on en
conclut $f_0=f_1$ dans $\mathrm{Hot}_{W}$, qed.
% archive sheet number: 263
